\subsubsection{Aplicação Cliente}
\label{sec:cliente}

A aplicação cliente é um \textit{Single Page Application} (SPA) \textit{mobile-first} que consome os 93 \textit{endpoints} da API e materializa, para os perfis de passageiro, motorista e administrador, os fluxos descritos na Seção~\ref{sec:interfaces}.

\paragraph{Processo de Desenvolvimento Assistido por IA}

A aplicação cliente foi integralmente desenvolvida com o apoio de ferramentas de inteligência artificial (IA) generativa, sob direção, revisão e validação do autor. Esse modelo de trabalho, frequentemente denominado \textit{AI-assisted development} ou programação assistida por IA, foi adotado de forma deliberada para acelerar a prototipação e a implementação, mantendo-se a autoria das decisões de requisitos, de arquitetura e de qualidade sob responsabilidade humana. Registra-se o uso dessas ferramentas em conformidade com as boas práticas de transparência acadêmica. O processo organizou-se em três etapas sequenciais e complementares, cada uma apoiada por uma ferramenta distinta.

Na primeira etapa, de geração inicial, a base do projeto foi produzida com a plataforma Lovable, ferramenta de geração de aplicações assistida por IA, que estabeleceu o esqueleto da aplicação, a \textit{stack} React com TanStack Router, Tailwind CSS e shadcn/ui, e um primeiro conjunto de telas funcionais. Na segunda etapa, de concepção visual, o Claude Design (da Anthropic) foi empregado para aprimorar as telas existentes e gerar novos protótipos de alta fidelidade, apoiando as decisões de \textit{layout}, de hierarquia visual e de sistema de design. Na terceira etapa, de implementação real, o assistente de programação agêntico Claude Code (também da Anthropic) foi utilizado para escrever, refatorar e evoluir o código de produção e para realizar a integração efetiva com o \textit{backend}, sempre a partir de instruções do autor e com revisão do resultado. É importante destacar que as telas geradas nas etapas iniciais foram posteriormente refatoradas e aprimoradas na fase de implementação: o código foi reorganizado na arquitetura de \textit{feature modules}, as chamadas à API foram centralizadas na camada de serviços, o controle de acesso por papel foi consolidado nos \textit{layouts} de rota e a apresentação foi unificada sob o sistema de design proprietário descrito a seguir. Dessa forma, a IA atuou como ferramenta de produtividade ao longo de todo o ciclo, da prototipação à integração, enquanto a estrutura, os padrões e a qualidade final do código resultam das decisões e da revisão conduzidas pelo autor.

O sistema de design da aplicação é proprietário: não se adotou nenhum \textit{design system} de terceiros (como Material Design ou Ant Design). Sua identidade visual, uma paleta de cores com variáveis semânticas (\textit{accent}, \textit{success}, \textit{danger}, \textit{warn}, entre outras), a tipografia (\texttt{Manrope} para texto e \texttt{JetBrains Mono} para dados numéricos) e a iconografia em traço, foi concebida na etapa de concepção visual, em um conjunto de primitivas escritas à mão acompanhado de um painel de ajuste em tempo real de paleta, raio de borda e densidade. Na etapa de implementação, essa identidade foi materializada na aplicação de produção sobre as primitivas acessíveis do shadcn/ui (Radix) e do Tailwind CSS: os \textit{tokens} de design foram portados para variáveis CSS e os componentes, reestilizados para a identidade do Movy. Trata-se, portanto, de uma identidade visual original implementada sobre primitivas de interface acessíveis, e não da adoção de um tema de terceiros.

\paragraph{Roteamento e Controle de Acesso}

O roteamento utiliza o mecanismo \textit{file-based} do TanStack Router, no qual a estrutura de nomes dos arquivos determina a árvore de URLs. \textit{Layouts} sem caminho (prefixados por sublinhado) não acrescentam segmento à URL, mas estabelecem nós sob os quais outras rotas são aninhadas, permitindo associar a um conjunto de rotas uma verificação de acesso comum. A autenticação é verificada uma única vez, no \textit{layout} \texttt{\_protected}, que envolve todas as rotas que exigem sessão; sobre essa base, dois \textit{layouts} adicionais, \texttt{\_admin} e \texttt{\_driver}, aplicam as verificações de papel, redirecionando usuários sem o papel correspondente. A composição resulta em uma cadeia de verificações (autenticado e administrador, por exemplo), sem que os segmentos de proteção apareçam na URL. Esse modelo é o reflexo, no cliente, do RBAC implementado pelos \textit{guards} do \textit{backend}.

\begin{figure}[H]
\centering
\setlength{\fboxsep}{8pt}
\fbox{%
\begin{minipage}{0.94\textwidth}
\small
\begin{tabularx}{\linewidth}{>{\bfseries}p{0.24\linewidth} X}
Público &
\texttt{/}, \texttt{/login}, \texttt{/signup}, \texttt{/public/trip-instances}, \texttt{/public/organizations}, \texttt{/public/plans}. \\
\texttt{\_protected} &
Rotas que exigem sessão: \texttt{/my-bookings}, \texttt{/profile}, \texttt{/trips/:orgId/:tripId/book}, \texttt{/setup}. \\
\texttt{\_admin} &
Subárvore administrativa: \texttt{/dashboard}, \texttt{/trips}, \texttt{/templates}, \texttt{/drivers}, \texttt{/vehicles}, \texttt{/organization}, \texttt{/financial}. \\
\texttt{\_driver} &
Subárvore operacional do motorista: \texttt{/my-trips} e acesso à tela compartilhada de detalhe da viagem. \\
\end{tabularx}
\end{minipage}}
\caption{Organização das rotas e dos \textit{guards} no cliente}
\label{fig:fe_rotas_guards}
\end{figure}

\paragraph{Padrão de Serviços e Hooks}

A Tabela~\ref{tab:fe_camadas} concretiza o fluxo unidirecional através das camadas, tomando como exemplo a funcionalidade ``minhas inscrições'': um serviço expõe operações de alto nível sobre os \textit{endpoints}; um \textit{hook} de caso de uso combina o serviço com estado e lógica de apresentação (busca, filtros), derivando o estado de carregamento da ausência simultânea de conteúdo e erro; e uma rota fina apenas decide entre os estados (carregando, erro, conteúdo) e delega a renderização.

\begin{table}[H]
\centering
\caption{Responsabilidades por camada no fluxo serviço $\rightarrow$ \textit{hook} $\rightarrow$ rota}
\label{tab:fe_camadas}
{\footnotesize
\setlength{\tabcolsep}{4pt}
\begin{tabularx}{\textwidth}{@{}>{\raggedright\arraybackslash}p{2.8cm} >{\raggedright\arraybackslash}p{4.7cm} X@{}}
\toprule
\textbf{Camada} & \textbf{Exemplo} & \textbf{Responsabilidade} \\
\midrule
Serviço &
\texttt{bookings.service.ts} &
Expõe operações de alto nível sobre os \textit{endpoints}, como listar e cancelar inscrições, escondendo os detalhes do transporte HTTP. \\
\addlinespace
\textit{Hook} de caso de uso &
\texttt{useBookings.ts} &
Invoca o serviço, mantém o estado da tela (conteúdo e erro) e deriva o estado de carregamento. \\
\addlinespace
Rota fina &
\texttt{my-bookings.tsx} &
Instancia o \textit{hook}, decide entre carregamento, erro e conteúdo, e delega a renderização aos componentes de apresentação. \\
\bottomrule
\end{tabularx}
}
\end{table}

\paragraph{Integração com a API}

A comunicação com o \textit{backend} é concentrada no cliente HTTP único (\texttt{lib/api.ts}) e organizada pela camada de serviços. Os aspectos mais relevantes dessa integração são:

Autenticação. Toda requisição é autenticada por padrão, com o cabeçalho \texttt{Authorization: Bearer}; chamadas a \textit{endpoints} públicos desabilitam esse comportamento explicitamente. A expiração do \textit{access token} é tratada de forma transparente pela renovação automática descrita na Seção~\ref{sec:autenticacao}.

Endpoints públicos. As telas do contexto público consomem rotas anônimas expostas pela API, \texttt{GET /public/trip-instances}, \texttt{GET /public/organizations} (e \texttt{/:slug}) e \texttt{GET /public/plans}, permitindo a descoberta de viagens, organizações e planos sem autenticação, em conformidade com o funil de conversão do produto.

Contexto multi-tenant. O isolamento por organização é responsabilidade do \textit{backend} (campo \texttt{organizationId} e \texttt{TenantFilterGuard}); no cliente, o contexto de papel deriva a organização administrada do usuário cruzando \texttt{GET /organizations/me} com \texttt{GET /memberships/me/role/:organizationId}, padrão também reutilizado para listar as organizações em que o usuário atua como motorista.

Apresentação temporal. A API armazena horários recorrentes em UTC; o cliente os converte e exibe sempre em horário de Brasília, e nunca expõe \textit{cron expressions} ou horários UTC literais ao usuário.
